% !TEX root = ../../ctfp-print.tex

\lettrine[lhang=0.17]{プ}{ログラマーたちはモナドをめぐる神話を} 作り上げてきた。モナドはプログラミングにおける最も抽象的で難解な概念の一つとされている。「理解できる」人と「できない」人がいる。多くの人にとって、モナドという概念を理解する瞬間は、まるで神秘的な体験のようだ。モナドはあまりにも多様な構成の本質を抽象化しているため、日常生活における良いアナロジーが存在しない。私たちは暗闇の中を手探りするしかない――まるで盲目の人々が象の異なる部分に触れては、「これはロープだ」「これは木の幹だ」「これはブリトーだ！」と得意げに叫ぶように。

誤解を正しておこう：モナドをめぐる神秘主義はすべて誤解の産物だ。モナドは非常にシンプルな概念だ。混乱の原因は、モナドの応用の多様性にある。

この記事のリサーチとして、ダクトテープ（別名：ダックテープ）とその応用を調べてみた。以下はそれで何ができるかのサンプルだ：

\begin{itemize}
  \tightlist
  \item
        ダクトの密封
  \item
        アポロ13号内のCO\textsubscript{2}スクラバーの修理
  \item
        いぼの治療
  \item
        AppleのiPhone 4の通話切断問題の修正
  \item
        プロムドレスの作成
  \item
        吊り橋の建設
\end{itemize}

\noindent
今、ダクトテープが何であるかを知らず、このリストから推測しようとしていると想像してほしい。幸運を祈る！

そこで、「モナドは〜のようなもの」という決まり文句のコレクションにもう一つ加えたい：モナドはダクトテープのようなものだ。その応用は多岐にわたるが、原理は非常にシンプルだ：物をくっつける。より正確には、物を合成する。

これは、多くのプログラマー、特に命令型プログラミングのバックグラウンドを持つプログラマーがモナドを理解するのに苦労する理由を部分的に説明している。問題は、私たちが関数合成の観点からプログラミングを考えることに慣れていないことだ。これは理解できる。私たちはしばしば、関数から関数へ直接渡す代わりに中間値に名前をつける。また、短いグルーコードのセグメントをヘルパー関数に抽象化する代わりにインライン展開する。Cでのベクトル長関数の命令型スタイルの実装を示す：

\begin{snip}{cpp}
double vlen(double * v) {
    double d = 0.0;
    int n;
    for (n = 0; n < 3; ++n)
        d += v[n] * v[n];
    return sqrt(d);
}
\end{snip}
これを、関数合成を明示的にした（スタイライズされた）Haskell版と比較してほしい：

\src{snippet01}
（ここで、さらに難解にするために、べき乗演算子\code{(\^{})}の二番目の引数を\code{2}に設定することで部分適用した。）

Haskellのポイントフリースタイルが常に優れていると主張しているのではない。ただ、関数合成がプログラミングにおけるすべての基礎にあるということだ。そして、私たちは実際に関数を合成しているにもかかわらず、Haskellはモナド的合成のために\code{do}記法と呼ばれる命令型スタイルの構文を提供するために多大な努力をしている。その使用法は後で見ることにする。しかしまず、なぜモナド的合成が必要なのかを説明しよう。

\section{クライスリ圏}

以前、通常の関数を装飾することで\hyperref[kleisli-categories]{ライターモナド}にたどり着いた。特定の装飾は、返り値を文字列と、より一般的にはモノイドの要素とペアリングすることによって行われた。このような装飾が関手であることを今や認識できる：

\src{snippet02}
続いて、装飾された関数――すなわちクライスリ射――を合成する方法を見つけた。クライスリ射は次の形の関数だ：

\src{snippet03}
ログの蓄積の実装は、合成の内部で行われた。

これで、クライスリ圏のより一般的な定義の準備が整った。圏$\cat{C}$と自己関手$m$から始める。対応する\newterm{クライスリ圏}（Kleisli category）$\cat{K}$は$\cat{C}$と同じ対象を持つが、射が異なる。$\cat{K}$における二つの対象$a$と$b$の間の射は、元の圏$\cat{C}$での射として実装される：
\[a \to m\ b\]
$\cat{K}$のクライスリ射を$a$と$m\ b$の間ではなく、$a$と$b$の間の射として扱うことを心に留めておくことが重要だ。

例では、$m$は何らかの固定されたモノイド\code{w}に対して\code{Writer w}に特殊化された。

クライスリ射が圏を形成するのは、それらに対して適切な合成を定義できる場合のみだ。結合的であり、すべての対象に対して恒等射を持つ合成が存在する場合、関手$m$は\newterm{モナド}（monad）と呼ばれ、結果として得られる圏はクライスリ圏と呼ばれる。

Haskellでは、クライスリ合成はフィッシュ演算子\code{>=>}を使って定義され、恒等射は\code{return}と呼ばれる多相関数だ。クライスリ合成を使ったモナドの定義を示す：

\src{snippet04}
モナドの定義には多くの同等な方法があり、これはHaskellエコシステムでの主要な定義ではないことを覚えておいてほしい。概念的なシンプルさと直感が気に入っているが、プログラミング時により便利な他の定義もある。それらについてはすぐに話す。

この定式化では、モナド則を非常に簡単に表現できる。Haskellでは強制できないが、等式推論に使用できる。これらは単にクライスリ圏の標準的な合成則だ：

\begin{snip}{haskell}
(f >=> g) >=> h = f >=> (g >=> h) -- 結合律
return >=> f = f                  -- 左単位律
f >=> return = f                  -- 右単位律
\end{snip}
この種の定義はモナドが本当に何であるかを表現している：装飾された関数を合成する方法だ。副作用や状態についてではなく、合成についてだ。後で見るように、装飾された関数はさまざまな効果や状態を表現するために使えるが、それはモナドの目的ではない。モナドは、装飾された関数の一端を別の装飾された関数の他端に結びつける粘着性のあるダクトテープだ。

\code{Writer}の例に戻ろう：ロギング関数（\code{Writer}関手のクライスリ射）は\code{Writer}がモナドであるため圏を形成する：

\src{snippet05}
\code{Writer w}のモナド則は\code{w}のモノイド則が満たされている限り満たされる（こちらもHaskellでは強制できない）。

\code{Writer}モナドには\code{tell}と呼ばれる便利なクライスリ射が定義されている。その唯一の目的は引数をログに追加することだ：

\src{snippet06}
後で他のモナド的関数のビルディングブロックとして使用する。

\section{フィッシュの解剖}

異なるモナドに対してフィッシュ演算子を実装すると、多くのコードが繰り返されており、容易に因数分解できることにすぐ気づく。まず、二つの関数のクライスリ合成は関数を返さなければならないため、型\code{a}の引数を取るラムダから始める実装になる：

\src{snippet07}
この引数でできることは\code{f}に渡すことだけだ：

\src{snippet08}
この時点で、型\code{m b}のオブジェクトと関数\code{g :: b -> m c}を持って、型\code{m c}の結果を生成しなければならない。これを行う関数を定義しよう。この関数は\newterm{バインド}（bind）と呼ばれ、通常は中置演算子の形で書かれる：

\src{snippet09}
すべてのモナドに対して、フィッシュ演算子を定義する代わりにバインドを定義できる。実際、Haskellの標準的なモナドの定義はバインドを使う：

\src{snippet10}
\code{Writer}モナドのバインドの定義を示す：

\src{snippet11}
確かに、フィッシュ演算子の定義よりも短い。

\code{m}が関手であるという事実を活用して、バインドをさらに分解することが可能だ。\code{fmap}を使って関数\code{a -> m b}を\code{m a}の中身に適用できる。これは\code{a}を\code{m b}に変換する。適用の結果は型\code{m (m b)}になる。これは求めているものとは少し違う――型\code{m b}の結果が必要だ――が、近づいている。必要なのは\code{m}の二重適用を潰す、あるいは平坦化する関数だ。そのような関数は\code{join}と呼ばれる：

\src{snippet12}
\code{join}を使って、バインドを次のように書き直せる：

\src{snippet13}
これによって、モナドを定義する第三の選択肢が得られる：

\src{snippet14}
ここでは、\code{m}が\code{Functor}であることを明示的に要求した。モナドの前の二つの定義ではそうする必要はなかった。なぜなら、フィッシュまたはバインド演算子をサポートする型コンストラクタ\code{m}は自動的に関手になるからだ。たとえば、バインドと\code{return}を使って\code{fmap}を定義することが可能だ：

\src{snippet15}
完全を期して、\code{Writer}モナドの\code{join}を示す：

\src{snippet16}

\section{\texttt{do}記法}

モナドを使ってコードを書く一つの方法は、クライスリ射を使って――フィッシュ演算子で合成して――作業することだ。このプログラミングモードはポイントフリースタイルの一般化だ。ポイントフリーコードはコンパクトで、しばしばかなりエレガントだ。ただし、一般的に理解しにくく、難解に近い場合がある。だから、ほとんどのプログラマーは関数の引数や中間値に名前をつけることを好む。

モナドを扱う場合、フィッシュ演算子よりバインド演算子を好むことを意味する。バインドはモナド値を受け取り、モナド値を返す。プログラマーはそれらの値に名前をつけることができる。しかし、それは大した改善ではない。本当に望んでいることは、それらをカプセル化するモナドコンテナではなく、通常の値を扱っているふりをすることだ。命令型コードはそのように機能している――グローバルログの更新などの副作用はほとんど見えない。そして、それがHaskellの\code{do}記法がエミュレートするものだ。

では、なぜモナドを使うのかと思うかもしれない。副作用を見えなくしたいなら、なぜ命令型言語を使わないのか？答えは、モナドが副作用をよりよく制御できるからだ。たとえば、\code{Writer}モナドのログは関数から関数へ渡され、グローバルに公開されることはない。ログを混乱させたりデータ競合を生じさせたりする可能性がない。また、モナド的コードはプログラムの残りの部分から明確に区分され、隔離されている。

\code{do}記法はモナド合成の単なる糖衣構文だ。表面上は命令型コードによく似ているが、バインドとラムダ式のシーケンスに直接変換される。

たとえば、\code{Writer}モナドのクライスリ射の合成を説明するために以前使った例を考えよう。現在の定義を使って、次のように書き直せる：

\src{snippet17}
この関数は入力文字列のすべての文字を大文字に変換し、単語に分割する。その間ずっとアクションのログを生成する。

\code{do}記法ではこのようになる：

\src{snippet18}
ここで、\code{upStr}は単なる\code{String}だ。\code{upCase}が\code{Writer}を生成するにもかかわらず：

\src{snippet19}
これは\code{do}ブロックがコンパイラによって次のように脱糖されるからだ：

\src{snippet20}
\code{upCase}のモナド的結果は\code{String}を取るラムダにバインドされる。この文字列の名前が\code{do}ブロックに現れる。行を読むとき：

\src{snippet21}
\code{upStr}が\code{upCase s}の結果を\emph{受け取る}と言う。

擬似命令型スタイルは、\code{toWords}をインライン展開するときにさらに顕著になる。\code{tell}の呼び出しで置き換える。これは文字列\code{"toWords "}をログに記録し、続いて\code{words}を使って文字列\code{upStr}を分割した結果で\code{return}を呼び出す。\code{words}は文字列を操作する通常の関数であることに注意してほしい。

\src{snippet22}
ここで、doブロックの各行は、脱糖されたコードで新しいネストされたバインドを導入する：

\src{snippet23}
\code{tell}がユニット値を生成するため、後続のラムダに渡す必要はない。モナド的結果の内容を無視すること（しかしその効果――ここではログへの貢献――は無視しない）は非常に一般的なので、その場合にバインドを置き換える特別な演算子がある：

\src{snippet24}
コードの実際の脱糖はこのようになる：

\src{snippet25}
一般に、\code{do}ブロックは行（またはサブブロック）で構成される。左矢印を使って残りのコードで利用可能な新しい名前を導入するか、副作用のためだけに実行される。コードの行の間にバインド演算子が暗黙的に存在する。ちなみに、Haskellでは\code{do}ブロックのフォーマットを中括弧とセミコロンで置き換えることができる。これにより、モナドをセミコロンのオーバーロードの方法として説明することが正当化される。

\code{do}記法を脱糖する際のラムダとバインド演算子のネストが、各行の結果に基づいて残りの\code{do}ブロックの実行に影響を与えるという効果があることに注意してほしい。このプロパティは複雑な制御構造を導入するために使用できる。たとえば、例外をシミュレートするためにだ。

興味深いことに、\code{do}記法に相当するものが命令型言語、特にC++でも応用されている。再開可能関数（リジューマブル関数）やコルーチンのことだ。C++の\urlref{https://bartoszmilewski.com/2014/02/26/c17-i-see-a-monad-in-your-future/}{futureがモナドを形成する}ことは秘密ではない。これは継続モナドの例であり、すぐに議論する。継続の問題は、合成が非常に困難なことだ。Haskellでは、\code{do}記法を使って「私のハンドラーがあなたのハンドラーを呼ぶ」というスパゲッティを、逐次コードに非常によく似たものに変換する。再開可能関数はC++でも同じ変換を可能にする。そして同じメカニズムを\urlref{https://bartoszmilewski.com/2014/04/21/getting-lazy-with-c/}{ネストされたループのスパゲッティ}をリスト内包表記や「ジェネレーター」に変換するために適用できる。ジェネレーターは本質的にリストモナドの\code{do}記法だ。モナドの統一的な抽象化なしでは、これらの問題はそれぞれ、言語へのカスタム拡張を提供することで対処されることが多い。Haskellでは、これらはすべてライブラリを通じて処理される。
